\section{Long Short-Term Memory}
\label{sec:lstm}

% TODO: fill this section once LSTM experiments are complete (see 7_lstm.ipynb)

\subsection{Architecture}

% TODO: describe the LSTM architecture.
%
% Key points to cover:
%   - Gating mechanism (input, forget, output gates)
%   - Cell state vs hidden state
%   - Number of layers, hidden size
%   - Dropout / regularization
%   - Output layer

\subsection{Training Setup}

% TODO: describe training procedure.

\subsection{Results}

% TODO: insert results table (MAE — full, calm, crisis).

% TODO: figure placeholder — predicted vs actual (LSTM)

% TODO: figure placeholder — training / validation loss curves

\subsection{Cross-Model Comparison}

% TODO: final summary table comparing all models on MAE (primary metric).
%
% Template:
%
% \begin{table}[H]
%   \centering
%   \caption{Summary comparison of all models (MAE on test set).}
%   \label{tab:final_comparison}
%   \begin{tabular}{lrrr}
%     \toprule
%     Model                       & MAE (full) & MAE (calm) & MAE (crisis) \\
%     \midrule
%     Linear Regression (raw)     & 0.638 & 0.488 & 1.987 \\
%     Linear Regression (eng.)    & 0.598 & ---   & ---   \\
%     GARCH(1,1)                  & 0.00583 & 0.00506 & 0.01281 \\
%     FNN                         & TODO  & TODO  & TODO  \\
%     RNN                         & TODO  & TODO  & TODO  \\
%     LSTM                        & TODO  & TODO  & TODO  \\
%     \bottomrule
%   \end{tabular}
% \end{table}

\subsection{Conclusion}

% TODO: summarise findings across all architectures.
%       Which model best handles the calm regime? The crisis regime?
%       What are the remaining open challenges (e.g., regime detection,
%       attention mechanisms, multivariate inputs)?
